\chapter*{Agradecimientos}
\addcontentsline{toc}{chapter}{Agradecimientos}
\textit{[Placeholder: Escribir aquí los agradecimientos a familiares, amigos, tutor (Prof. Walter Hernández) y la Ilustre Universidad Central de Venezuela]}

\chapter*{Resumen}
\addcontentsline{toc}{chapter}{Resumen}

El sistema PortalAsig, desarrollado para la gesti\'{o}n acad\'{e}mica de la Facultad de
Ciencias de la Universidad Central de Venezuela, opera desde el a\~{n}o 2016 sobre
una arquitectura monol\'{i}tica que acumul\'{o} seis a\~{n}os de deuda t\'{e}cnica sin
proceso sistem\'{a}tico de mantenimiento. Esta situaci\'{o}n, agravada por un contexto
operativo sin equipo de desarrollo de tiempo completo y con colaboradores
acad\'{e}micos temporales, gener\'{o} un grado de acoplamiento que dificultaba cualquier
intervenci\'{o}n y hac\'{i}a depender el sistema del conocimiento t\'{a}cito de cada
mantenedor.

El presente trabajo documenta el redise\~{n}o del backend de PortalAsig hacia una
arquitectura de microservicios que descompone el sistema en dominios acotados e
independientes ---UAA (gesti\'{o}n de identidad), Site (sitios acad\'{e}micos), Notify
(notificaciones) y PAIMA (autogesti\'{o}n del mantenimiento)--- cada uno con su
propia base de datos, su esquema de versionado y su ciclo de despliegue aut\'{o}nomo.
El frontend heredado, desarrollado en Vue.js~2, fue adaptado para operar como
cliente distribuido: se reestructur\'{o} el store Vuex en m\'{o}dulos por dominio, se
implementaron guards de navegaci\'{o}n con siete niveles de permisos, se migr\'{o} la
autenticaci\'{o}n a OAuth2/JWT con cookies y refresh token, y se configuraron tres
instancias de Axios que reflejan la topolog\'{i}a de los microservicios. Sobre esta
arquitectura se implementaron t\'{e}cnicas sistem\'{a}ticas de mantenimiento de software
que abarcan todo el ciclo de vida del c\'{o}digo: integraci\'{o}n continua con validaci\'{o}n
automatizada, an\'{a}lisis est\'{a}tico de c\'{o}digo, pruebas de integraci\'{o}n sobre bases de
datos reales, pruebas end-to-end y medici\'{o}n de cobertura. Se instrumentaron
adem\'{a}s mecanismos de observabilidad distribuida que recolectan m\'{e}tricas,
agregan registros de ejecuci\'{o}n y visualizan correlaciones entre eventos del
ecosistema.

Sobre esta infraestructura sensorial se construy\'{o} el servicio PAIMA (PortalAsig
Intelligent MAPE-K), que ejecuta un ciclo de autogesti\'{o}n del mantenimiento
mediante el modelo MAPE-K, donde las fases de An\'{a}lisis y Planificaci\'{o}n se
delegan a un modelo de inteligencia artificial. PAIMA monitorea el ecosistema,
detecta anomal\'{i}as, diagnostica causas, clasifica incidencias, notifica al
administrador y propone acciones de remediaci\'{o}n que el operador humano eval\'{u}a y
aprueba.

La validaci\'{o}n mediante pruebas de integraci\'{o}n, validaci\'{o}n end-to-end y
simulaci\'{o}n de caos demostr\'{o} que el ecosistema a\'{\i}sala fallos, mantiene m\'{e}tricas
objetivas de calidad y reduce la carga cognitiva del mantenedor. El resultado es
un sistema desacoplado, observable y autogestionado que puede ser mantenido por
colaboradores acad\'{e}micos temporales sin depender del conocimiento t\'{a}cito.

\vspace{1em}
\noindent\textbf{Palabras clave:} arquitectura de microservicios, mantenimiento de
software, integraci\'{o}n continua, observabilidad distribuida, computaci\'{o}n auton\'{o}mica,
MAPE-K, inteligencia artificial, autogesti\'{o}n del mantenimiento.

\newpage
\tableofcontents
\newpage
% \listoffigures  % TODO: habilitar cuando todas las figuras sean reales
\newpage
